% chapters/05-conclusion.tex — บทที่ 5 สรุปและข้อเสนอแนะ
\chapter{สรุปและข้อเสนอแนะ}
\label{ch:conclusion}

\section{สรุปผลการศึกษา}

งานวิจัยนี้ได้พัฒนาระบบอัตโนมัติสำหรับการคำนวณต้นทุนต่อหน่วยของวัสดุก่อสร้าง
จากข้อมูลของกระทรวงพาณิชย์และสถานประกอบการค้าปลีกสมัยใหม่ทั่วประเทศ
สำเร็จตามวัตถุประสงค์ที่ตั้งไว้ทั้ง 3 ข้อ ดังนี้

\begin{enumerate}
  \item \textbf{ระบบดึงข้อมูลอัตโนมัติ:} พัฒนาระบบที่รองรับ 11 แหล่งข้อมูล
        (3 ภาครัฐ + 8 Modern Trade) ผ่าน 4 กลยุทธ์ ได้แก่ PDF parsing,
        JSON API, ScrapingBee proxy + Browser Rendering, และ Manual Upload
  \item \textbf{การคำนวณต้นทุนต่อหน่วย:} รองรับการคำนวณ BoQ ครบ 5 ประเภทงาน
        (\texttt{wall\_tile}, \texttt{column\_beam}, \texttt{rebar},
        \texttt{paint}, \texttt{brick}) ตามหลักเกณฑ์การคำนวณราคากลาง
        ตุลาคม 2560 พร้อมเชื่อมโยงราคาจากหลายแหล่งและบันทึกผลใน
        \texttt{fact\_calculations} dimension
  \item \textbf{การวิเคราะห์เชิงเวลาและพื้นที่:} ระบบให้บริการ Linear Trend
        Analysis, Percentage Change Analysis, Outlier Detection (z-score),
        และ Regional Aggregator (median ราคา 6 ภาค)
\end{enumerate}

\section{การมีส่วนร่วมเชิงวิชาการ (Contribution)}

จากการทบทวนวรรณกรรม พบว่างานวิจัยก่อนหน้า \autocite{rao2015,chunyaem2019,liu2024}
ได้พัฒนาแนวทางเฉพาะส่วน เช่น Web Scraping, AI classification, หรือ Price
Forecasting แต่ยังขาดงานวิจัยที่บูรณาการกระบวนการทั้งหมดเข้าด้วยกัน. งานวิจัย
นี้ \emph{เติมเต็มช่องว่างดังกล่าว} โดยเชื่อมโยง pipeline ตั้งแต่การดึงข้อมูล
จัดการมาตรฐาน คำนวณต้นทุน ไปจนถึงการวิเคราะห์เชิงสถิติในระบบเดียว ที่
deploy บน Cloudflare Workers edge network.

\section{ข้อจำกัดของการศึกษา}

\begin{enumerate}
  \item \textbf{Bot challenge:} 5 จาก 8 retail sources (Thai Watsadu, Global
        House, BnB, SCG Home, DoHome) ถูก block ด้วย Cloudflare bot challenge
        หรือเป็น SPA ที่โหลดราคาด้วย XHR หลัง hydrate
  \item \textbf{Data freshness:} ราคาประเมินทางการ (CGD) เปลี่ยนรายไตรมาส
        ระบบใช้ TTL 30 วัน + cron weekly snapshot ซึ่งเพียงพอสำหรับ research
        scale แต่อาจไม่ทันต่อตลาดในช่วงราคาผันผวน
  \item \textbf{Sample coverage:} ระบบทดสอบจริงเฉพาะ 24 รายการวัสดุ ครอบคลุม
        77 จังหวัด แต่ regional coverage ยังบาง (ส่วนใหญ่ default ที่ province 10
        กรุงเทพฯ); งานเชิงพื้นที่จึงใช้ AI-researched seed prices ใน 5 ภาค
\end{enumerate}

\section{ข้อเสนอแนะสำหรับงานวิจัยในอนาคต}

\begin{enumerate}
  \item \textbf{Reverse-engineer SPA APIs:} วิเคราะห์ XHR endpoint ของ
        DoHome, SCG Home, Global House เพื่อเข้าถึงข้อมูลโดยไม่ต้อง render
        full browser
  \item \textbf{Forecasting model:} ขยายจาก linear trend ไปสู่ ARIMA หรือ
        Prophet เพื่อพยากรณ์ราคาในอนาคต
  \item \textbf{Trust weighting:} กำหนดน้ำหนักของแหล่งข้อมูลแบบ adaptive
        (Government 2x Retail) เมื่อรวมเป็น index
  \item \textbf{Province expansion:} ขยายการเก็บข้อมูล regional ให้ครบ
        77 จังหวัด โดยใช้ TPSO regional bulletin เป็น primary source
  \item \textbf{Public API:} เปิด \texttt{/api-docs} ให้ผู้ประมาณราคา
        third-party ใช้ฟรี พร้อม rate-limit policy
\end{enumerate}

\section{บทสรุปสำคัญ}

ระบบที่พัฒนาขึ้นพิสูจน์ว่าเป็นไปได้ในทางปฏิบัติที่จะ \emph{รวมศูนย์} ข้อมูล
ราคาวัสดุก่อสร้างจากหลายแหล่งให้กลายเป็น \emph{single source of truth}
สำหรับการประมาณราคา BoQ. การพัฒนา ระบบบน edge platform (Cloudflare Workers)
ทำให้ระบบมีต้นทุนการดำเนินงานต่ำ (scale-to-zero) และ latency ระดับ <50ms
ทั่วโลก ซึ่งเหมาะสมกับการใช้งานในงานก่อสร้างภาครัฐและภาคเอกชน. งานวิจัยนี้
จึงมีศักยภาพในการนำไปใช้สนับสนุนการบริหารงบประมาณภาครัฐให้มีความโปร่งใสและ
สอดคล้องกับสภาพตลาดจริงมากยิ่งขึ้น.
